\documentclass[11pt]{article}

\usepackage[margin=1in]{geometry}
\usepackage{amsmath,amssymb,amsthm,mathtools}
\usepackage{enumitem}
\usepackage{booktabs}
\usepackage{tabularx}
\usepackage{microtype}
\usepackage[hidelinks]{hyperref}
\usepackage{xurl}

\newtheorem{definition}{Definition}
\newtheorem{remark}{Remark}
\newtheorem{lemma}{Lemma}

\newcommand{\TV}{\mathrm{TV}}

\title{Digest-Bound Replay Plans and State Declarations:\\
\large Equivalence, Minimality, and Change Control (Series Synthesis)}
\author{}
\date{Unified archive note v0.19 (2026-03-23; archive v667)}

\begin{document}
\maketitle

\begin{abstract}
The archive publishes brief, referee-grade anonymity papers while withholding large artifact bundles (logs, scripts, certificates) unless requested.
This is only credible if (i) the public receipt surface commits to \emph{which} replay workflow would validate the claim, and (ii) any long-lived system state assumptions are pinned in a comparable way.

This note standardizes two digest-bound adjuncts already used by the worked example:
\texttt{plan\_spec\_id} (a canonical hash of the replay-plan specification) and \texttt{state\_decl\_id} (a canonical hash of the state-surface declaration).
We then give a small, operational taxonomy for drift and equivalence:
which changes are ``publication-only'' (re-run the verifier), and which changes are ``claim-affecting'' (recertify).
\end{abstract}

\paragraph{Later-terse-paper citation posture.}
Later terse papers should cite this note (Synthesis~18) only when the live issue is the digest-bound semantic core of \texttt{plan\_spec\_id} / \texttt{state\_decl\_id}, the P0/P1/P2 or S0/S1/S2 drift class itself, or the compact first-reopen-owner default that decides which narrower or wider owner regains authority first once one such digest-bound shortcut has gone stale.
They should stop at Synthesis~12 for receipt tuple semantics, at Synthesis~19 or Synthesis~21 when the live issue is the exact horizon-accounting regime or adaptive per-step certificate semantics, at Synthesis~20 for replay checklist / replay-verdict mechanics, at Synthesis~30 for deployed-surface ownership, at State~1 or Certified~C when the live issue is the stateful attenuation argument or recertification policy owned there, and at Synthesis~32 together with Synthesis~34 and Synthesis~74--75 once the question has already become release-stage routing, outward-facing wrapper status, or whether a later terse paper may stop, widen, or reopen.

\paragraph{Scope.}
This is an interface/change-control note.
It does \emph{not} propose new privacy mechanisms.
It exists so other papers can stay tight:
when a paper says ``replay plan $\pi$'' or ``state contract $\sigma$,'' it should mean a digest-bound object, not prose.

\paragraph{Imports.}
Identifier conventions are centralized in Synthesis~8.\cite{series:notation}

Contract-object vocabulary is imported from Synthesis~0.\cite{series:contractobjects}
Receipt line-item semantics are imported from Synthesis~12.\cite{series:receiptitems}
Trace-interface and threat-window identities are imported from Synthesis~3/4.\cite{series:traceifaces,series:threatwindows}
The motivating end-to-end instantiation is the worked example (Synthesis~17).\cite{series:workedexample}
State-dependent accounting motivation is imported from State~1.\cite{series:state1}
Recertification policy under drift is imported from Certified~C.\cite{series:recert}
Deployed-surface ownership (state vs. contact surface vs. ENF/projection) is imported from Synthesis~30.\cite{series:deployedsurfaces}

\paragraph{Exports.}
(i) Minimal semantic cores for \texttt{plan\_spec} and \texttt{state\_decl},
(ii) digest-binding rules (\texttt{plan\_spec\_id}, \texttt{state\_decl\_id}),
(iii) a drift/equivalence taxonomy that tells deployers when to republish, when to re-run verifiers, and when to recertify,
(iv) a compact digest-bound first-reopen-owner card saying which narrower or wider owner regains authority first once one plan/state shortcut has gone stale,
(v) an optional publication-obligation adjunct that says which small public outputs should be refreshed once the action class is known, and
(vi) an optional release-spine adjunct that names the stable successor-maintenance stages once the narrower maintenance objects exist.

\section{Replay plans: what a receipt promises}
Many papers export a numeric budget, but the public object that makes the budget meaningful is a \emph{receipt line item}.\cite{series:receiptitems}
A line item cannot ship the full artifact bundle in a source-only release, so it must at least commit to:\
(a) \emph{which} replay workflow would validate the claim, and
(b) \emph{what} evidence that workflow expects.

\begin{definition}[Replay plan specification]
A \emph{replay plan specification} \texttt{plan\_spec} is a short, deterministic description of the replay workflow referenced by \texttt{replay\_hook}.
It should contain only semantics a third party must know to (i) request the right private artifacts and (ii) deterministically recompute the public outputs.
\end{definition}

\paragraph{Minimal semantic core.}
To minimize needless churn while remaining referee-grade, we recommend that a \texttt{plan\_spec} include exactly:
\begin{enumerate}[leftmargin=*]
\item \textbf{Owner and scope:} the paper/module that owns the replay semantics; the claim fields the plan is intended to validate (typically \texttt{exposure\_nf\_id}, \texttt{tw\_id}, and optional \texttt{state\_decl\_id}).
\item \textbf{Inputs:} a list of required artifacts by name \emph{and} a short input-contract statement (e.g., ``trace bundle includes raw transcripts + eligibility filter outputs + randomness commitments'').
\item \textbf{Determinism:} the intended interpreter/toolchain version or a verifier binary digest.
\item \textbf{Checks:} a small list of named predicates the verifier applies (e.g., conformance, sample-selection, accountant recomputation).
\item \textbf{Outputs:} the expected public outputs (e.g., \texttt{verifier\_report.json} schema hash; receipt acceptance bit; derived OINL fields).
\end{enumerate}
Avoid timestamps, hostnames, and other non-semantic metadata.
If deployments want to ship operator notes, they should be separate, non-hashed fields.

\begin{definition}[Digest-binding for replay plans]
Define
\[\texttt{plan\_spec\_id} := \texttt{sha256}(\texttt{canon}(\texttt{plan\_spec})),\]
where \texttt{canon} is a deterministic canonicalization (sorted keys, no whitespace, stable number formatting).
A receipt line item that includes \path{replay_hook.plan_id} should also include \path{replay_hook.plan_spec_id} so ``same plan id, different replay'' cannot hide off-surface.\cite{series:receiptitems,series:workedexample}
\begin{remark}[Plan certificates and transparency logs are replay evidence]
A \texttt{plan\_spec\_id} can also bind operational commitments that live outside the core anonymity math: for example,
a schedule-only deployment can publish a signed plan certificate and an inclusion proof in a plan transparency log (PTL) and treat those objects as named inputs to the replay workflow.
Operational~B provides a concrete instantiation of this pattern.\cite{series:operationalB}
\end{remark}

\end{definition}

\begin{remark}[Why this does not overreach]
\texttt{plan\_spec\_id} does \emph{not} certify the privacy claim.
It certifies the \emph{replay promise}: what an auditor will do if they request the evidence.
This is the minimum needed for non-equivocation and mechanical diffing of receipts across releases.
\end{remark}

\section{State declarations: what state a bound assumes}
Many anonymity claims are implicitly stateful (caches, reputation, long-lived monitors, adaptive fallback gates).
When a bound depends on such state, a human label like \texttt{state\_contract\_id} is not a stable interface:
it can drift without changing the label.

\begin{definition}[State-surface declaration]
A \emph{state declaration} \texttt{state\_decl} is a signed adjunct that commits to which long-lived control-plane state surfaces are in scope for the claim.
It should be treated as part of the witness interface for stateful deployments.\cite{series:state1,series:workedexample}
\end{definition}

\paragraph{Minimal semantic core.}
A \texttt{state\_decl} should include:
\begin{enumerate}[leftmargin=*]
\item \textbf{Surface inventory:} a list of named state surfaces (e.g., ``fallback usage counters,'', ``cache eviction policy,'' ``monitoring summaries'') with explicit include/exclude flags.
\item \textbf{Retention and reset:} the reset unit (epoch, window) and retention horizon for each included surface.
\item \textbf{Read/write actors:} which roles can write the surface (router, monitor, operator) and what is considered public leakage vs. internal state.
\item \textbf{Binding to the claim window:} the intended \texttt{tw\_id} (or family) the declaration is meant to accompany.
\end{enumerate}
As with plans, avoid non-semantic metadata.

\begin{definition}[Digest-binding for state surfaces]
Define
\[\texttt{state\_decl\_id} := \texttt{sha256}(\texttt{canon}(\texttt{state\_decl\_core})),\]
where \texttt{state\_decl\_core} is the minimal semantic core and \texttt{canon} is a deterministic canonicalization.
Receipts for stateful claims should carry \texttt{state\_decl\_id}.\cite{series:receiptitems,series:workedexample}
\end{definition}

\begin{lemma}[Conditioning on a published state witness bounds joint leakage]\label{lem:tv_state_witness}
Let $x,x'$ be two secrets (e.g., two lookup targets) and let $(W,Y)$ be the pair of
(i) a \emph{state witness} $W=g(S)$ derived from internal state $S$ and
(ii) a published view $Y$ (receipt line items, exposure-normal-form projections, or replay outputs).
Write $P_x$ for the joint law of $(W,Y)$ under secret $x$ and $P_{x'}$ analogously.
Then
\[
\TV(P_x, P_{x'}) \le \TV(P_x^W, P_{x'}^W) \;+\; \sup_{w}\,\TV\!\big(P_x^{Y\mid W=w},\,P_{x'}^{Y\mid W=w}\big).
\]
In particular, if $W$ is public and independent of the secret (e.g., a public refresh timer or a public ``bucket-staleness'' class),
then the joint distinguishing advantage of $(W,Y)$ is controlled by the worst-case conditional leakage of $Y$ given $W$.
\end{lemma}

\begin{proof}
For each $w,y$,
\[
P_x(w,y)-P_{x'}(w,y)=\big(P_x(w)-P_{x'}(w)\big)P_x(y\mid w)+P_{x'}(w)\big(P_x(y\mid w)-P_{x'}(y\mid w)\big).
\]
Summing absolute values over $(w,y)$ and using $\sum_y P_x(y\mid w)=1$ gives
$\|P_x-P_{x'}\|_1 \le \|P_x^W-P_{x'}^W\|_1 + \sup_w \|P_x^{Y\mid w}-P_{x'}^{Y\mid w}\|_1$.
Divide by $2$ to obtain the claim.
\end{proof}

\begin{remark}[How to use this lemma in the series]
State \emph{declarations} pin which state surfaces are in scope.
When a proof actually conditions on a public state summary (rate-limit class, bucket-staleness class, refresh epoch),
publishing that summary as an explicit receipt line item (or as part of the exposure normal form) makes the dependence auditable,
and Lemma~\ref{lem:tv_state_witness} turns ``statefulness'' into a conditional-budget accounting step.
See the state series for examples of when omitting such witnesses breaks naive horizon composition.\cite{series:state1}
\end{remark}


\section{Equivalence and drift: an operational taxonomy}
This section is deliberately operational: it tells a deployer what action is required when an id changes.

\begin{table}[t]
\centering
\small
\begin{tabularx}{\linewidth}{@{}lX@{}}
\toprule
Change & Minimal required action \\
\midrule
\textbf{ENF-ID or TW-ID changes} & New claim object (new receipt line item and ABOM/OINL); any reuse must be explicit via lineage/coarsening.\cite{series:traceifaces,series:threatwindows}\\
\textbf{StateDecl changes} (\texttt{state\_decl\_id} changes) & Treat as claim-affecting for stateful bounds: republish receipt/ABOM; if the new surface could increase leakage, trigger recertification.\cite{series:state1,series:recert}\\
\textbf{PlanSpec changes} (\texttt{plan\_spec\_id} changes) & Always republish the receipt/ABOM to pin the new replay promise. Whether the privacy claim must be recertified depends on whether the verifier/accountant semantics changed (taxonomy below).\\
\bottomrule
\end{tabularx}
\caption{Digest-bound ids turn ``is this the same claim?'' into a mechanical diff.}
\label{tab:diff}
\end{table}

\subsection*{Plan-spec drift classes (P0/P1/P2)}
Let \texttt{plan\_spec} change from \texttt{plan\_spec\_id} from $h_0$ to $h_1$.
Classify the change by what an auditor would compute on the same artifact bundle:
\begin{itemize}[leftmargin=*]
\item \textbf{P0 (non-semantic):} the verifier output and acceptance predicate are unchanged (e.g., refactoring, documentation, performance).
\emph{Action:} republish the new \texttt{plan\_spec\_id}; no new privacy evidence is required.
\item \textbf{P1 (checker-strengthening / determinism fixes):} acceptance may change only by rejecting previously-accepted bundles (more checks, tighter conformance, improved determinism).
\emph{Action:} republish and re-run the verifier on the old artifact bundle; if it still accepts, the claim can be reused.
\item \textbf{P2 (claim-affecting semantics):} the accountant, sampling rule, or acceptance predicate changes in a way that could accept bundles that would previously be rejected, or could change the computed bound.
\emph{Action:} treat as recertification-relevant drift (Certified~C): mint a new claim or run incremental/targeted audits as appropriate.\cite{series:recert}
\end{itemize}
This classification is intentionally conservative: without publishing full artifacts, the public guarantee is only as stable as the published replay semantics.

\subsection*{State-decl drift classes (S0/S1/S2)}
Similarly, classify \texttt{state\_decl} drift by whether the new declaration expands what can depend on the secret:
\begin{itemize}[leftmargin=*]
\item \textbf{S0 (metadata-only):} renaming or rewording with no surface change.
\item \textbf{S1 (surface tightening):} removing a surface from scope or reducing retention.
\item \textbf{S2 (surface expansion):} adding a surface to scope, increasing retention, or allowing new actors to write secret-dependent state.
\end{itemize}
\emph{Action:} S1 can often reuse the old bound conservatively; S2 should be treated as claim-affecting unless the owning state paper supplies an explicit attenuation lemma.

\subsection*{Optional adjunct: release-action matrices}
To keep mechanism and worked-example papers from rephrasing the same drift prose, a release may publish a tiny digest-bound \emph{release-action matrix}: a field-family map from observed diffs to the \emph{minimal default action}. This adjunct is optional but useful whenever users are expected to compare receipts across releases without reading the whole owner stack.

\begin{table}[t]
\centering
\small
\begin{tabularx}{\linewidth}{@{}>{\raggedright\arraybackslash}p{0.24\linewidth} >{\raggedright\arraybackslash}p{0.18\linewidth} X@{}}
\toprule
Field family & Default action & Escalation / ownership note \\
\midrule
Publication metadata, log anchors, watch text & Refresh only & Publication-layer drift only; keep the certified interface fixed. \\
\texttt{plan\_spec\_id}, replay-plan catalog digests & Replay or P0/P1/P2 handling & Use the P0/P1/P2 taxonomy above; P2 remains recertification-relevant. \\
\texttt{contact\_surface\_id}, line-item knobs, \texttt{budget\_value}, \texttt{evidence\_id} & Replay changed surfaces & Same interface, new effective service surface or evidence. If reader-visible transcript semantics change, escalate by minting a new \texttt{exposure\_nf\_id}. \\
\texttt{tw\_id}, \texttt{state\_decl\_id}, required \texttt{exposure\_nf\_id}s, primary separation class & Full recertification & Threat-window, state-surface, projection, or primary-guard drift changes the certified interface or its guard. \\
\bottomrule
\end{tabularx}
\caption{Compact field-family-to-action defaults. Synthesis~30 owns the state/contact/ENF ownership split that prevents contact-surface drift from being mistaken for ENF drift.}
\label{tab:releaseactionmatrix}
\end{table}

\begin{remark}[Why this reduces redundancy]
A release-action matrix is not a new receipt primitive.
It is a compact, digest-bound rendering of the P0/P1/P2 and S0/S1/S2 rules together with the state/contact/ENF ownership split.
The worked example (Synthesis~17) now ships one concrete matrix so users can tell whether a diff means ``refresh'', ``rerun the published replay plans'', or ``stop claiming continuity and recertify'' without every paper restating the same three paragraphs.\cite{series:workedexample,series:deployedsurfaces}
\end{remark}

\paragraph{A compact digest-bound first-reopen-owner card.}
A release-action matrix answers ``what class is this diff?'' but later terse papers often need one narrower routing answer: \emph{which already-owned public note should speak first once that class is known?}
This note therefore exports one tiny digest-bound \emph{first-reopen-owner card}: a field-family map from drift class to the first already-owned public owner together with the adjacent widen-only-if condition.
The card is routing-only.
It does not create a new theorem, replay packet, or release wrapper.
It simply keeps digest-bound drift from being flattened into vague ``something changed; see the stack'' prose once later terse papers start reusing the same replay/state packet under space pressure.

\begin{table}[t]
\centering
\small
\begin{tabularx}{\linewidth}{@{}>{\raggedright\arraybackslash}p{0.22\linewidth} >{\raggedright\arraybackslash}p{0.24\linewidth} X@{}}
\toprule
Drift family after classification & First reopen owner & Why stop there first / widen only if \\
\midrule
\texttt{plan\_spec\_id} under P0/P1, with the same certified theorem packet and the same published budget row & Replay-local owner (for the current synthesis stack: Synthesis~20; for the current Boss Fight public packet: Boss Fight~C) & The replay promise changed while the certifiable claim row is still the same. Widen to evaluator notarization only if attempt-prefix, randomness binding, or retry discipline is now live. \\
\texttt{plan\_spec\_id} under P2, or any replay change that can alter the computed bound or acceptance semantics & Claim-root / recertification owner (for the current stack: the owning mechanism paper or Certified~C) & The verifier semantics themselves may have changed what can be certified, so replay-local wrappers are already too late. \\
Same replay/state packet, but the live issue is now the exact horizon regime, horizon-total rule, or adaptive per-step certificate semantics & Accounting / conditional-certificate owner (for the current synthesis stack: Synthesis~19 or Synthesis~21) & The digest-bound packet is still the same, but the honest moving part is no longer replay mechanics as such: it is the exact horizon-accounting or adaptive certificate object that later terse papers are trying to reuse. \\
\texttt{contact\_surface\_id}, declared knob vectors, or line-item evidence hooks, with stable \texttt{exposure\_nf\_id} and \texttt{tw\_id} & Deployed-surface / knob-surface owner (for the current stack: Synthesis~30 or Boss Fight~B) & The declared contacted-reader surface or declared knob contract changed while the certified interface stayed fixed. Widen to replay only if checker contract rather than declared knob surface is now live. \\
\texttt{state\_decl\_id} under S1/S2, with stable \texttt{exposure\_nf\_id} and \texttt{tw\_id} & State-local owner (for the current stack: State~1 or the narrower state specialization) & In-scope secret-dependent state, retention, or writer roles changed. Widen to recertification only if no imported attenuation lemma still covers the expansion conservatively. \\
\texttt{exposure\_nf\_id}, \texttt{tw\_id}, or primary-separation-class drift & Claim root / recertification owner & The certified interface or its guard changed, so no downstream replay or wrapper note should speak first. \\
Wrapper-only release-stage or public-status drift, after the narrower theorem/state/replay owners are already fixed & Release-stage / public-status owner (for the current stack: Synthesis~32 or Synthesis~34; for legacy public wrappers only: Release~A) & The live issue is now the successor-maintenance packet or outward-facing wrapper rather than the narrower digest-bound semantic owner. \\
\bottomrule
\end{tabularx}
\caption{A compact digest-bound first-reopen-owner card. Unlike the worked example's fully concrete card, this one stays at the drift-class level and only names the first already-owned public stop once one maintained plan/state shortcut has gone stale.}
\label{tab:firstreopenowner}
\end{table}

\paragraph{One tiny worked-bundle first-reopen-owner drill.}
Use the maintained worked-example compare packet rather than a hypothetical checker swap.
The shipped compare report \path{example_compare_report.json} (\texttt{report\_id=worked-example-compare-report-v1}) classifies the canned base$\to$successor diff as \texttt{replay\_changed\_surfaces}, keeps the continuity decision at \texttt{same\_certified\_interface\_replayed\_surfaces}, and names the exact replay scope: changed line items \texttt{fallback\_contact\_surface} and \texttt{fallback\_tiered\_vector} with replay plans \texttt{eval-fallback-cct-v1} and \texttt{eval-fallback-vector-v1}.
The companion action matrix \path{example_release_action_matrix.json} (\texttt{action\_matrix\_id=worked-example-release-action-matrix-v1}) matches that diff to the replay-required rules rather than to interface/state recertification, while the compare report's stable guard fields keep \texttt{state\_decl\_id=sha256:ddfbb55199a6ecdb225cbafc60a74c234ae3334fb40528fde0a2ec2ce6de3f20} and \texttt{tw\_id=sha256:c1ff81fc84049a2160a21487f61a4565e03c5263f1437649c4299d1354b35c18} fixed.
So the first honest public stop is replay-local: in the current synthesis stack that means Synthesis~20, not Certified~C and not the release wrapper.
If that same successor instead changed \texttt{state\_decl\_id} away from the published value above, the same card would fail closed to the state / recertification owners (State~1 and Certified~C).
If all narrower semantic rows stayed fixed and only the outward-facing successor notice or closure wrapper changed, the honest widening stop would move to Synthesis~32 or Synthesis~34 rather than replay or state prose.\cite{series:workedexample,series:state1,series:recert,series:receiptaccounting,series:conditionalbudget,series:statusenvelopes,series:releasespine,series:auditorreplay}

\paragraph{Optional adjunct: publication-obligation profiles.}
A release-action matrix answers ``what class is this diff?'' but not necessarily ``what public outputs should be refreshed now?'' For that second question, a release may publish a tiny digest-bound \emph{publication-obligation profile}: a map from action class (refresh only / replay changed surfaces / full recertification) to the minimal public output roles that should appear in the successor publication package. Typical roles are: fresh release binding, fresh compare report, fresh replay verdict, fresh successor continuity verdict, or a Certified~C style recertification packet. This keeps change classification separate from release-management duties while still making both checkable.

\begin{remark}[Why split classification from obligations]
The same action class can recur across many papers and releases, but maintainers still need one small object saying what the publication package must contain. Publishing that profile once lets mechanism papers and release notes cite a checked object instead of improvising ad hoc checklists. The worked example now ships one concrete obligation profile alongside its action matrix and successor verdict.\cite{series:workedexample,series:auditorreplay}
\end{remark}

\paragraph{Optional adjunct: successor continuity verdicts.}
When a release compares one specific successor against one specific base receipt, it may publish a tiny digest-bound \emph{successor continuity verdict}: a capsule that points to the compare profile/report, records the matched release-action rule and replay scope, and states the continuity decision (for example, \path{same_certified_interface_replayed_surfaces} versus ``stop continuity and recertify''). This is still not a new receipt primitive. It is just the portable outcome object that lets later papers cite one base-to-successor conclusion instead of re-narrating the same diff, replay, and continuity prose. The worked example now ships one such verdict adjunct, while Synthesis~20 owns the replay-verdict fields it draws from.\cite{series:workedexample,series:auditorreplay}

\paragraph{Optional adjunct: publication-closure verdicts.}
An obligation profile says what roles a successor package should contain; a publication-closure verdict says whether one concrete base$\to$successor package actually satisfied that obligation. This tiny digest-bound adjunct simply binds each required publication role to the concrete shipped artifact path (or records that a role is missing). It is useful precisely because it is narrower than change classification and narrower than replay: later papers can cite one checkable package-completeness capsule instead of restating release-package contents in prose. The worked example now ships one such closure verdict for its canned replay-required Case~B path.\cite{series:workedexample,series:auditorreplay}

\paragraph{Optional adjunct: public-status wrappers and order objects.}
A successor continuity verdict says whether continuity holds, a publication-closure verdict says whether the required outputs were actually shipped, and a successor-derivation graph says what depended on what.
When later papers mainly need the outward-facing answer, a release may publish one tiny digest-bound successor-lineage notice to collate those narrower objects into a single public-status capsule.
When later papers mainly need the stable maintenance order rather than the concrete DAG, a release may publish one tiny digest-bound release spine naming the canonical stages diff $\to$ classify $\to$ obligate $\to$ close $\to$ summarize $\to$ support $\to$ cite $\to$ quote $\to$ clause.
The worked example now ships both objects, so later notes can cite one wrapper or one order object rather than re-narrating the whole maintenance stack.\cite{series:workedexample,series:auditorreplay}

\paragraph{Optional adjunct: downstream normal forms.}
Even with those objects in hand, later papers may still drift when they compress the public-status sentence into smaller claims, citations, quotes, and reusable clauses.
For that narrower downstream problem, this series now uses one compact home: Synthesis~31.
It treats claim-support maps, citation maps, quote maps, clause packs, and the tiny normal-form wrapper as one small auditable suffix rather than a growing list of unrelated micro-adjuncts.
That keeps this note about drift classes, equivalence, and publication duties rather than sentence choreography.\cite{series:downstreamnormalforms,series:workedexample}

\section{Checklist for keeping papers non-redundant}
A mechanism paper should not define replay/state semantics in prose.
Instead it should:
\begin{enumerate}[leftmargin=*]
\item Export a receipt line item with \texttt{exposure\_nf\_id}, \texttt{tw\_id}, and, when relevant, \texttt{state\_decl\_id}.
\item Point replay mechanics at \texttt{plan\_id} \emph{and} \texttt{plan\_spec\_id}.
\item When describing drift/updates, cite this note for the P0/P1/P2 and S0/S1/S2 rules rather than restating them.
\item If cross-release user diffing matters, publish a tiny release-action matrix instead of re-explaining field-family actions in prose.
\item If release-management duties would otherwise be repeated, publish a tiny publication-obligation profile that maps action classes to the refreshed public outputs.
\item If package completeness for one concrete successor matters, publish a tiny publication-closure verdict that binds each required role to one shipped artifact (or marks it missing).
\item If later papers mainly need the outward-facing status, publish a tiny successor-lineage notice that cites the compare report plus the continuity/closure verdicts instead of rephrasing them.
\item If later papers also need the maintenance order itself, publish a tiny successor-derivation graph instead of restating the compare/replay/continuity/closure dependency chain.
\item If later papers need one stable stage vocabulary for release evolution, publish a tiny release spine instead of re-narrating the diff/classify/obligate/close/summarize/support/cite/quote/clause chain in prose.
\item If later papers need sentence-level support bindings, publish a tiny successor-claim support map instead of re-listing which narrower object backs each clause of the public status sentence.
\item If later papers keep re-explaining what the smallest sufficient citation is for each maintenance question, publish a tiny successor-citation map instead of improvising citation ladders in prose.
\item If later papers keep paraphrasing slightly different fields from the same maintenance object, publish a tiny successor-quote map that names the exact field or field-list worth quoting.
\item If later papers still keep rebuilding the same short wrapper clauses after field choice is fixed, publish a tiny successor-clause pack that reuses those checked fields while preserving the declared cite path.
\end{enumerate}

\begin{thebibliography}{99}\small

\bibitem{series:contractobjects}
Anonymous.
\newblock Contract Objects for Anonymous-DHT Privacy Budgets and Claims.
\newblock Synthesis~0, The-One-True-Archive (2026).

\bibitem{series:receiptitems}
Anonymous.
\newblock Receipt Line Items for Anonymity Claims: A Minimal Tuple Schema for Published Budgets.
\newblock Synthesis~12, The-One-True-Archive (2026).

\bibitem{series:traceifaces}
Anonymous.
\newblock Trace Interfaces for Anonymous DHTs.
\newblock Synthesis~3, The-One-True-Archive (2026).

\bibitem{series:threatwindows}
Anonymous.
\newblock Threat Models and Repetition Windows for Anonymous DHTs.
\newblock Synthesis~4, The-One-True-Archive (2026).

\bibitem{series:workedexample}
Anonymous.
\newblock Worked Example: Receipt/Interlock and Drift Comparison.
\newblock Synthesis~17, The-One-True-Archive (2026).

\bibitem{series:operationalB}
Anonymous.
\newblock Auditable Trace Privacy with Abort: Plan-Pinning and Drift Receipts for Side-Channel Resistant Transport.
\newblock Operational series Paper~B, 2026. (This archive.)

\bibitem{series:bossfightB}
Anonymous.
\newblock AnonDHT in the Boss Fight Regime: Compiling a MaxL Budget into Lookup Knobs.
\newblock Boss Fight series Paper~B, The-One-True-Archive (2026).

\bibitem{series:bossfightC}
Anonymous.
\newblock Proof-Carrying Budgets for Anonymous DHTs.
\newblock Boss Fight series Paper~C, The-One-True-Archive (2026).

\bibitem{series:evalnotary}
Anonymous.
\newblock Notarized Anonymity Certificates Under Retries.
\newblock Evaluation series Paper~1, The-One-True-Archive (2026).

\bibitem{series:abomoinl}
Anonymous.
\newblock ABOM and OINL for Anonymity Claims.
\newblock Anonymity series Paper~C, The-One-True-Archive (2026).

\bibitem{series:releaseproperty}
Anonymous.
\newblock Release Properties, Lineage Tags, and Ledger Receipts for Anonymous Deployments.
\newblock Release and destination series Paper~A, The-One-True-Archive (2026).

\bibitem{series:state1}
Anonymous.
\newblock State-Dependent Anonymity Budgets for Anonymous DHTs.
\newblock AnonDHT State series Paper~1, The-One-True-Archive (2026).

\bibitem{series:recert}
Anonymous.
\newblock Disagreement-Gated Recertification and Boundary Audits for Anonymous DHT Routers.
\newblock Certified series Paper~C, The-One-True-Archive (2026).

\bibitem{series:notation}
Anonymous.
\newblock Notation and Minimal Model for the One-True-Archive.
\newblock Series synthesis note (Synthesis~8), The-One-True-Archive (2026).

\bibitem{series:releasespine}
Anonymous.
\newblock Release-Stage Spine for Public Evidence Maintenance.
\newblock Series synthesis note (Synthesis~32), 2026. (This archive.)

\bibitem{series:statusenvelopes}
Anonymous.
\newblock Public Status Envelopes for Release Maintenance.
\newblock Series synthesis note (Synthesis~34), 2026. (This archive.)

\bibitem{series:receiptaccounting}
Anonymous.
\newblock Receipt Horizon Accounting and Regime Packets.
\newblock Series synthesis note (Synthesis~19), 2026. (This archive.)

\bibitem{series:conditionalbudget}
Anonymous.
\newblock Conditional Budget Certificates for Adaptive Replay Horizons.
\newblock Series synthesis note (Synthesis~21), 2026. (This archive.)

\bibitem{series:deployedsurfaces}
Anonymous.
\newblock Deployed Observation Surfaces for Anonymous DHTs: Control-plane knobs in IPFS/libp2p delegated routing and OHTTP evolution.
\newblock Series synthesis note (Synthesis~30), The-One-True-Archive (2026).

\bibitem{series:downstreamnormalforms}
Anonymous.
\newblock Downstream Normal Forms for Successor Maintenance: Summary, Support, Citation, Quote, and Clause Discipline.
\newblock Series synthesis note (Synthesis~31), 2026. (This archive.)

\bibitem{series:auditorreplay}
Anonymous.
\newblock Auditor Replay for Anonymous-DHT Receipt Line Items.
\newblock Series synthesis note (Synthesis~20), The-One-True-Archive (2026).

\end{thebibliography}

\end{document}
